\documentclass[10pt,a4paper]{amsart}
\usepackage[margin=19mm]{geometry}
\usepackage{amsmath,amssymb,mathtools}
\usepackage[scr=rsfs,cal=euler]{mathalfa}
\usepackage{xfrac}
\usepackage[unicode,hidelinks,bookmarks=false]{hyperref}
\newcommand{\ie}{i.e.\ }
\newcommand{\cf}{cf.\ }
\newcommand{\resp}{respectively\ }
\newcommand{\Subsection}{\subsection}
\providecommand{\nobreakdash}{\nobreak\hbox{-}}
\setlength{\emergencystretch}{2em}
\multlinegap0pt
\usepackage{amssymb}
\usepackage{mathtools}
\usepackage{xcolor}
\usepackage[shortlabels]{enumitem}
\usepackage{algorithm}\usepackage{algpseudocode}
\usepackage{tcolorbox}
\newtheorem{lemma}{Lemma}
\title{Evolving Local Corrections for Global Constructions in Combinatorics}
\author{Gergely B\'erczi}
\date{Source: arXiv:2603.06692v1 (CC BY 4.0)}
\begin{document}
\maketitle

\noindent\emph{Source and licence.} Evolving Local Corrections for Global Constructions in Combinatorics. Version arXiv:2603.06692v1 (CC BY 4.0), distributed by arXiv under the Creative Commons Attribution 4.0 International licence, http://creativecommons.org/licenses/by/4.0/, as recorded in the arXiv licence metadata for this exact version and shown on the abstract page https://arxiv.org/abs/2603.06692v1. Full licence notice: \texttt{licenses/arxiv-2603.06692-CC-BY-4.0.txt}.\par\smallskip

\noindent\textit{Editorial scope.} Counted unit: the article's lemma lem:profile (source0.tex lines 397--406). The article's complete original proof of it is delivered with it (source0.tex lines 408--415, one \texttt{proof} environment of 8 lines). No mathematics is written, repaired or re-derived: the statement and the proof are the article's own lines, in the article's own notation, and no retained line is substituted or re-flowed.  No further source-local context is needed: every cross-reference of every retained line resolves inside the two delivered spans.  Nothing was omitted from any retained span, so the package carries no omission record.  No retained line contains a citation, so the wrapper carries no bibliography and no bibliography entry.  No retained line contains a figure environment, an image-inclusion command or drawing code, so no image is delivered and the package declares no asset dependency.  Editorial formatting only, in addition: an AMS \texttt{amsart} preamble that restates the article's own declarations and macros used by the retained lines, namely the environments \texttt{lemma} on separate counters, together with the standard packages those lines need.  The retained environments print in this document as Lemma 1 in order of appearance, whereas the article prints the same items with the numbering of its own full document.  The complete native source of the article as distributed in the arXiv e-print for this exact version is bundled unchanged as \texttt{sources/195887/source0.tex}.
\begin{lemma}[Recovering neighbor-degree profiles from one deletion]\label{lem:profile}
Let $G=(U\sqcup V,E)$ be bipartite. Fix $u\in U$ and let $G' = G-u$.
Let $g_k$ be the number of vertices $v\in V$ with $\deg_G(v)=k$, and let $c_k$ be the number of vertices $v\in V$ with $\deg_{G'}(v)=k$.
Let $a_k$ be the number of neighbors $v\in N(u)$ with $\deg_G(v)=k$.
Then for all $k\ge 0$,
\[
c_k = g_k - a_k + a_{k+1}.
\]
Consequently, knowing the sequences $(g_k)_{k\ge 0}$ and $(c_k)_{k\ge 0}$ determines $(a_k)_{k\ge 0}$ uniquely by downward recursion.
\end{lemma}

\begin{proof}
Deleting $u$ decreases the degree of a vertex $v\in V$ by $1$ iff $v\in N(u)$, and leaves it unchanged otherwise.
Hence, among vertices in $V$ of degree $k$ in $G$, exactly $a_k$ drop to degree $k-1$ in $G'$. This contributes $-a_k$ to the count at degree $k$ in $G'$.
Meanwhile, among degree-$(k+1)$ vertices in $G$, exactly $a_{k+1}$ drop into degree $k$ in $G'$, contributing $+a_{k+1}$.
All other vertices retain their degree. Thus $c_k = g_k - a_k + a_{k+1}$.
If degrees are bounded by $\Delta$, set $a_{\Delta+1}=0$ and solve backwards:
$a_\Delta = g_\Delta - c_\Delta$, then $a_{\Delta-1}= g_{\Delta-1}+a_\Delta - c_{\Delta-1}$, etc.
\end{proof}

\end{document}
